\clearpage\thispagestyle{empty}
\pdfbookmark[0]{Quellen zu Beobachtungen und Elemente des Kometen (1. August 1844)}{cometaug1844observations-sources}
\begin{center}\textbf{Quellen zu Beobachtungen und Elemente des Kometen (1. August 1844)}\end{center}
\noindent Diese KI-Übertragung ins heutige Deutsch folgt dem Bericht auf S.~466--467 der \href{https://gdz.sub.uni-goettingen.de/id/PPN236007467}{\textit{Werke}, Band~VI} (Göttingen, 1874). Der nicht unterzeichnete Beitrag in den \textit{Astronomischen Nachrichten}, Band~XXII, Nr.~511, \href{https://api.digitale-sammlungen.de/iiif/image/v2/bsb10049447_00064/full/1800,/0/default.jpg}{Sp.~111--112}, ausgegeben am 1. August 1844, berichtet über Beobachtungen von Carl Friedrich Gauß. Der Katalog der Bayerischen Staatsbibliothek datiert den gebundenen Band auf 1845.
\par\smallskip\noindent Der Erstdruck enthält außerdem Beobachtungen von Rümker und die von Rümker und Funk berechneten Elemente. Der Auszug in der Werkausgabe endet mit der Anzahl von Gauß’ Vergleichen. Die Tabelle folgt der Werkausgabe bei den Zeitzeichen, dem Dezimalpunkt und den beiden Kometenzeichen. KI-Übertragung und KI-Satz.
\clearpage
